\documentclass[a4paper,12pt]{article}
\usepackage[utf8]{inputenc}
\usepackage[spanish]{babel}
\usepackage[T1]{fontenc}
\usepackage{graphicx}
\usepackage[most]{tcolorbox}
\usepackage{geometry}
\usepackage{tabularx}
\usepackage{enumitem}
\usepackage{helvet}
\usepackage{setspace}
\usepackage{array}

\geometry{top=2cm, bottom=2.5cm, left=3cm, right=2cm}
\renewcommand{\familydefault}{\sfdefault}
\setstretch{1.15}

\definecolor{bannergreen}{RGB}{160, 210, 60}

\renewcommand{\arraystretch}{1.5}
\setlength{\arrayrulewidth}{0.6pt}

% --- DATOS ESPECÍFICOS ---
\newcommand{\distrito}{4117}
\newcommand{\escuela}{Escuela Técnica N° X}
\newcommand{\docente}{Ferreyra Gustavo David}
\newcommand{\materia}{Mecánica de los Fluidos}
\newcommand{\curso}{5to Año}
\newcommand{\division}{A}
\newcommand{\ciclo}{2026}

\newcommand{\makeheader}[1]{%
    \begin{center}
        \begin{minipage}{0.25\textwidth}
            \centering
            \includegraphics[width=2.2cm]{distritos/4117/templates/logo1.jpg}
        \end{minipage}
        \hfill
        \begin{minipage}{0.25\textwidth}
            \centering
            \includegraphics[width=2.2cm]{distritos/4117/templates/logo2.jpg}
        \end{minipage}
        
        \vspace{0.3cm}
        \begin{tcolorbox}[colback=bannergreen, colframe=bannergreen, arc=20pt, center, halign=center, fontupper=\bfseries\normalsize, boxrule=0pt, width=0.98\textwidth, top=1pt, bottom=1pt]
            #1
        \end{tcolorbox}
    \end{center}
    \vspace{0.4cm}
}

\begin{document}

\makeheader{PLANIFICACIÓN ANUAL \ciclo}

\begin{tabularx}{\textwidth}{|l|X|l|X|}
    \hline
    \textbf{Escuela:} & \escuela & \textbf{Curso:} & \curso \division \\ \hline
    \textbf{Docente:} & \docente & \textbf{Materia:} & \materia \\ \hline
    \textbf{Distrito:} & \distrito & \textbf{Ciclo:} & \ciclo \\ \hline
\end{tabularx}

\vspace{0.8cm}

\section*{Contenidos}
\begin{itemize}[label=\textbullet, leftmargin=1cm]
    \item Principio de Pascal
    \item Principio de Arquímedes
    \item Ecuaciones de Continuidad
    \item Teorema de Bernoulli
    \item Número de Reynolds
\end{itemize}

% --- OBJETIVOS ---
\section*{Objetivos}
\begin{itemize}[label=\textbullet]
    \item \textbf{General:} Comprender y aplicar los principios fundamentales de la estática y dinámica de fluidos para la resolución de problemas de ingeniería.
    \item \textbf{Específicos:}
    \begin{itemize}
        \item Aplicar los principios de Pascal y Arquímedes en el análisis de sistemas hidráulicos.
        \item Analizar el flujo de fluidos mediante las ecuaciones de continuidad y el Teorema de Bernoulli.
        \item Diferenciar regímenes de flujo (laminar y turbulento) mediante el cálculo del Número de Reynolds.
        \item Resolver problemas prácticos de transporte de fluidos y presión.
    \end{itemize}
\end{itemize}

% --- METODOLOGÍA ---
\section*{Metodología}
\begin{itemize}[label=\textbullet]
    \item \textbf{Enfoque:} Teórico-Práctico. Se prioriza la resolución de problemas complejos y la aplicación de modelos matemáticos a fenómenos físicos.
    \item \textbf{Estrategias:}
    \begin{itemize}
        \item Clases expositivas con soporte visual y deducción de fórmulas.
        \item Resolución de guías de trabajos prácticos enfocadas en la hidráulica.
        \item Análisis de casos reales de flujo en tuberías y conductos.
        \item Experimentación guiada sobre los principios fundamentales de los fluidos.
    \end{itemize}
\end{itemize}

% --- EVALUACIÓN ---
\section*{Evaluación}
\begin{itemize}[label=\textbullet]
    \item \textbf{Criterios:}
    \begin{itemize}
        \item Precisión en la aplicación de fórmulas y resolución de problemas numéricos.
        \item Capacidad de análisis crítico de los fenómenos de flujo observados.
        \item Cumplimiento y entrega en tiempo y forma de los trabajos prácticos.
        \item Participación activa en la resolución de ejercicios en clase.
    \end{itemize}
    \item \textbf{Instrumentos:}
    \begin{itemize}
        \item Evaluaciones escritas parciales.
        \item Entregas de trabajos prácticos corregidos.
        \item Instancias de defensa oral de los ejercicios resueltos.
    \end{itemize}
\end{itemize}

\end{document}